\section{Related Work}

% Historically, the bispectrum emerged from the study of the higher-order statistics of non-Gaussian random processes, and was used as a tool to measure the non-Gaussianity of a signal \cite{tukey1953spectral, brillinger1991some, nikias1993signal}. The work of Yellott and Iverson \cite{yellott1992uniqueness} established that every integrable function with compact support is completely identified\textemdash up to translation\textemdash by its three-point autocorrelation and thus its bispectrum (Appendix \ref{appendix:autocorrelation}). Later work of Kakarala tied the bispectrum to the more general formulation of Fourier analysis and thus defined it for all compact groups \cite{kakarala1993group, kakarala2012bispectrum}. Kakarala went on to prove that it is complete \cite{kakarala2009completeness} in both commutative and non-commutative formulations (Appendix \ref{appendix:completeness}).

\textbf{The Triple Correlation.} The triple correlation has a long history in signal processing \cite{tukey1953spectral, brillinger1991some, nikias1993signal}. It originally emerged from the study of the higher-order statistics of non-Gaussian random processes, but its invariance properties with respect to translation have been leveraged in texture statistics~\cite{yellott1993implications} and
data analysis in neuroscience~\cite{deshpande2023third}, as well as early multi-layer perceptron architectures in the 1990's~\cite{delopoulos1994invariant,kollias1996multiresolution}. The triple correlation was extended to groups beyond translations in \cite{kakarala1992triple}, and its completeness with respect to general compact groups was established in \cite{kakarala2009completeness}. To the best of our knowledge, the triple correlation has not previously been introduced as a method for achieving invariance in convolutional networks for either translation or more general groups.


% As our proposed method can be used as a variant of a pooling mechanism, we first review existing pooling that have been employed for CNNs in the literature, and that have been, or not extended to $G$-CNNs.

\textbf{Pooling in CNNs.} Pooling in CNNs typically has the dual objective of coarse graining and achieving local invariance. While invariance is one desiderata for the pooling mechanism, the machinery of group theory is rarely employed in the computation of the invariant map itself. As noted in the introduction, max and average pooling are by far the most common methods employed in CNNs and $G$-CNNs. However, some approaches beyond strict max and average pooling have been explored.
%Pooling can serves other various purposes, such as down-scaling (coarse-graining), increasing the receptive field of convolutional filters across layers, and reducing computational complexity and memory usage.
% Consequently, numerous approaches have been proposed in academic literature, extending  methods, and their mixed-max-average pooling variants~\cite{lee2015generalizing}.
Soft-pooling addresses the lack of smoothness of the max function and uses instead a smooth approximation of it, with methods including polynomial pooling~\cite{wei2019building} and learned-norm \cite{gulcehre2014learned}, among many others \cite{estrach2014signal,eom2020alphaintegration, shahriari2017multipartite,shi2016rank,bieder2021comparison,stergiou2021refining,czaja2021maximal,kumar2018ordinal}.
% p-pooling \cite{estrach2014signal}, alpha-integration \cite{eom2020alphaintegration}, multi-partite~\cite{shahriari2017multipartite}, rank-based \cite{shi2016rank}, Smooth-Maximum \cite{bieder2021comparison}, Soft pooling \cite{stergiou2021refining}, Maxfun pooling~\cite{czaja2021maximal}, Ordinal pooling~\cite{kumar2018ordinal}.
Stochastic pooling~\cite{zeiler2013stochastic,} reduces overfitting in CNNs by introducing randomness in the pooling, yielding mixed-pooling~\cite{yu2014mixed}, max pooling dropout~\cite{wu2015max}, among others~\cite{tong2016hybrid,zhai2017s3pool,graham2014fractional}
% Hybrid pooling~\cite{tong2016hybrid},
%  Stochastic Spatial Sampling (S3 pooling)~\cite{zhai2017s3pool} and Fractional Max pooling~\cite{graham2014fractional}.


\textbf{Geometrically-Aware Pooling.} Some approaches have been adopted to encode spatial or structural information about the feature maps, including spatial pyramid pooling~\cite{he2015spatial}, part-based pooling~\cite{zhang2012pose}, geometric $L_p$ pooling \cite{feng2011geometric} or pooling regions defined as concentric circles \cite{qi2018concentric}. In all of these cases, the pooling computation is still defined by a max. These geometric pooling approaches are reminiscent of the Max $G$-Pooling for $G$-CNNs introduced by ~\cite{cohen2016group} and defined in Section \ref{sec:g-pooling}, without the explicit use of group theory.

\textbf{Higher-Order Pooling.} Average pooling computes first-order statistics (the mean) by pooling from each channel separately and does not account for the interaction between different feature maps coming from different channels. Thus, second-order pooling mechanisms have been proposed to consider correlations between features across channels~\cite{lin2015bilinear,gao2019global}, but higher-orders are not investigated. Our approach computes a third-order polynomial invariant; however, it looks for higher-order correlations within the group rather than across channels and thus treats channels separately. In principle, these approaches could be combined.

% \paragraph{Universality and Invariance in Neural Networks} Researchers have explored the universality of $G$-invariant neural networks and showed that neural networks with $G$-invariant linear layers can approximate any $G$-invariant function~\cite{maron2019universality} as long as higher-order tensors are available. This is a different question than the one we tackle here, which focuses on exploiting the smallest $G$-invariant map that is not lossy. In \cite{yarotsky2022universal}, universality of $G$-invariant neural networks is also studied, and qualified as ``complete'' in the sense of approximation theory, which is different from the meaning of the word ``complete,'' in the sense of invariant theory used here.
